\documentclass[conference]{IEEEtran}
% EE302 PS-3 report: battery monitoring and passive/active cell balancing.
% Compile with pdfLaTeX (e.g. upload this folder to Overleaf). Figures live in figures/.
% A figure file that is missing is drawn as a labelled box, so the document always compiles.
% Red [..] marks are items that still need real data or names.

\usepackage[T1]{fontenc}
\usepackage{amsmath,amssymb}
\usepackage{graphicx}
\usepackage{booktabs}
\usepackage{xcolor}
\usepackage{cite}
\usepackage{url}
\usepackage[hidelinks]{hyperref}
\graphicspath{{figures/}}

\newcommand{\tbd}[1]{\textcolor{red}{\textbf{[#1]}}}
% \figslot{file}{width}{placeholder height}{description}
\newcommand{\figslot}[4]{%
  \IfFileExists{figures/#1}{\includegraphics[width=#2]{#1}}{%
    \fbox{\parbox[c][#3][c]{\dimexpr#2-2\fboxsep-2\fboxrule\relax}{\centering\small\textcolor{gray}{%
    \textbf{Image placeholder}\\[2pt] #4\\[2pt] \texttt{\detokenize{figures/#1}}}}}}}

\begin{document}

\title{Battery Monitoring and Passive and Active\\ Cell Balancing of a Four-Cell Li-ion Pack:\\ Simulation and Arduino Hardware}

\author{
\IEEEauthorblockN{Vidhi Chandak}
\IEEEauthorblockA{B24356\\IIT Mandi}
\and
\IEEEauthorblockN{Vitthal Seth}
\IEEEauthorblockA{B24357\\IIT Mandi}
\and
\IEEEauthorblockN{Shashwat Saha}
\IEEEauthorblockA{B23292\\IIT Mandi}
\and
\IEEEauthorblockN{Rishi Bansal}
\IEEEauthorblockA{B23288\\IIT Mandi}
}

\maketitle

\begin{abstract}
We design, simulate and build a battery management system (BMS) for four series 18650 Li-ion cells on an Arduino
Mega 2560 with a 16-bit ADS1115 converter and four DS18B20 sensors. One controller provides monitoring, passive
balancing with relay-switched 47\,$\Omega$ resistors, active balancing with a relay-switched capacitor ladder,
latched fault handling, alerts and a manual stop, selected over a serial link. A Simscape model with a verified
averaged switched-capacitor stage predicted that both methods reduce a 107\,mV spread to about 11\,mV, passive
balancing in 1.55\,h and active balancing in 1.05\,h. On hardware, four runs were recorded. In a 36\,min passive run the controller bled the correct cells at about 85\,mA, removing 34 and 30\,mAh from the
two highest cells; their rested voltages fell by 14.1 and 23.1\,mV, against 41\,mV expected in total from the
removed charge, with 0.31\,Wh dissipated and no measurable cell heating. In three
active runs (2.5\,h in total) no transfer was measurable: the relays shuttle 470\,$\mu$F capacitors about once
per second, which moves 20--60\,$\mu$A per pair. The fault logic caught an open sense wire, and the warning logic
caught later intermittent contacts. Calibration against a multimeter reduced the channel error from up to 48\,mV
(41\,mV RMS) to $\pm1$\,mV. From the measured voltages we estimate 6.7\,h for passive balancing of the
pack and over 26\,000\,h for the present active stage, reduced to 24--196\,h by MOSFET switching at 1\,kHz. Code, raw data, schematics and graphs:
\url{https://github.com/highdreameater/EE-302-Control-sys-proj}.
\end{abstract}

\begin{IEEEkeywords}
battery management system, cell balancing, switched-capacitor equaliser, Li-ion, Arduino, fault detection
\end{IEEEkeywords}

\section{Introduction}
In a series pack, differences in capacity, internal resistance and self-discharge make the cells drift apart in
state of charge (SOC). The pack is then limited by its weakest cell: discharge stops when the lowest cell reaches
cut-off and charge stops when the highest cell reaches its limit. \emph{Passive} balancing dissipates the excess
charge of the high cells in resistors. \emph{Active} balancing moves charge between cells with capacitors or
inductors~\cite{daowd2011,gallardo2014}. Switched-capacitor equalisers are attractive for small packs because they
need no magnetic components~\cite{gallardo2014,lim2025}.

Following Problem Statement~3, the objectives are to (i) monitor the voltage and temperature of a four-cell pack
and calibrate the measurement, (ii) implement monitoring, passive, active, fault and stop modes with alerts in one
controller, (iii) implement resistor-based passive and switched-capacitor active balancing, (iv) compare the
methods in simulation and on hardware, and (v) estimate how long each method needs to balance the pack.

\begin{figure}[t]
\centering
\figslot{block_diagram.png}{\columnwidth}{4.2cm}{System block diagram}
\caption{System overview. The PC selects the mode over USB serial and records the log.}
\label{fig:block}
\end{figure}

\section{Hardware Design}

\subsection{Controller, sensing and pack}
The pack is four 18650 cells (about 2600\,mAh, 4.2\,V full, 2.8\,V cut-off) in series at about 15.7--15.9\,V. An
Arduino Mega 2560 runs the controller every 1\,s and logs one 22-column CSV row per second at 115\,200\,baud
(Figs.~\ref{fig:block} and~\ref{fig:photo}). Table~\ref{tab:io} lists the channels.

Each stack node $j$ reaches the ADS1115~\cite{ads1115} through a divider of $R_\text{top}=33$\,k$\Omega$ and
$R_\text{bot}=11$\,k$\Omega$,
\begin{equation}
V_{\text{node},j} = k\,V_{\text{ADC},j}, \qquad k=\frac{R_\text{top}+R_\text{bot}}{R_\text{bot}}=4.0,
\label{eq:div}
\end{equation}
and the cell voltages follow from one single-ended and three differential conversions,
\begin{equation}
\begin{aligned}
v_1 &= k V_{A0}, & v_2 &= k\,(V_{A1}-V_{A0}),\\
v_3 &= k\,(V_{A3}-V_{A1}), & v_4 &= k\,(V_{A2}-V_{A3}).
\end{aligned}
\label{eq:cells}
\end{equation}
At $\pm2.048$\,V full scale the resolution is 62.5\,$\mu$V, or 0.25\,mV per cell. Readings are averaged over four
conversions. Each cell has a DS18B20 (12-bit, 0.0625\,$^\circ$C) on one 1-Wire bus; conversions are started at the
end of a period and read at the next, so the 750\,ms conversion does not delay the loop.

\begin{table}[t]
\caption{I/O channels and scalings}
\label{tab:io}
\centering
\small
\begin{tabular}{@{}lll@{}}
\toprule
Signal & Pin / channel & Note \\
\midrule
Cell 1 & A0 (single-ended) & $\times4.0$ \\
Cell 2 & A1 $-$ A0 & $\times4.0$ \\
Cell 3 & A3 $-$ A1 & $\times4.0$ \\
Cell 4 & A2 $-$ A3 & $\times4.0$ \\
Temperatures $T_1$--$T_4$ & D2, 1-Wire & 0.0625\,$^\circ$C/LSB \\
Bleed relays K1--K4 & D22--D25 & active-LOW \\
Capacitor relays K5--K8 & D26--D29 & active-LOW \\
Fault LED / buzzer & D13 / D8 & solid = fault \\
Alert LED & D9 & on for any alert \\
Stop push-button & D3 & to GND, pull-up \\
\bottomrule
\end{tabular}
\end{table}

\begin{figure}[t]
\centering
\figslot{photo_hardware.jpg}{\columnwidth}{5cm}{Photo of the assembled board}
\caption{Assembled hardware.}
\label{fig:photo}
\end{figure}

\subsection{Passive balancing stage}
Each cell has a 47\,$\Omega$, 2\,W resistor switched by a relay (Fig.~\ref{fig:schpass}). At 4.0\,V
\begin{equation}
I_b=\frac{v_i}{R_b}\approx85\ \text{mA},\qquad P_b=\frac{v_i^2}{R_b}\approx0.34\ \text{W}.
\label{eq:bleed}
\end{equation}

\begin{figure}[t]
\centering
\figslot{schematic_passive.png}{\columnwidth}{6cm}{Measurement and passive balancing schematic}
\caption{Measurement dividers and relay-switched bleed resistors.}
\label{fig:schpass}
\end{figure}

\subsection{Active balancing stage}
The active stage is an adjacent-cell switched-capacitor ladder~\cite{pascual1997}: three 470\,$\mu$F, 50\,V
capacitors, each with a 4.7\,$\Omega$ series resistor, switched by four SPDT relays (Fig.~\ref{fig:schact}).
Released, the capacitors sit across cells 1--3; energised, across cells 2--4. Each relay acts as a half-bridge
across one cell. One shuttle is 500\,ms in each position.

\begin{figure}[t]
\centering
\figslot{schematic_active.png}{\columnwidth}{6cm}{Relay-switched capacitor ladder schematic}
\caption{Relay-switched capacitor ladder.}
\label{fig:schact}
\end{figure}

\subsection{Safety and indication}
Every output is switched off and the state is latched on cell over-voltage (4.35\,V), under-voltage (2.80\,V),
over-temperature (55\,$^\circ$C), two consecutive failed temperature reads, the serial \texttt{stop} command or the
push-button. A latched fault is cleared only by \texttt{reset} once all readings are back in range. Alerts (status
WARN) are raised for a cell below 3.30\,V in any mode, a temperature above 45\,$^\circ$C, a rested reading that
jumps by more than 30\,mV (a loose sense wire), and, in active mode, neighbouring cells more than 300\,mV apart;
balancing pauses during an alert. A fault gives a solid LED and a continuous buzzer; an alert gives a blinking LED,
a separate alert light and, for low voltage, a short chirp every 5\,s. No cell was charged during the trials.
A capacitor test emulator was not used. The protection and mode logic was instead verified in the Simscape model
(manual stop, over-temperature and low-voltage cases) and in a host simulation of the firmware (low-voltage alert,
stop, reset and mode changes) before the pack was connected, and the fault detection was then exercised on the
pack itself (Section~\ref{sec:hw}).

\section{Modelling and Simulation}
\label{sec:sim}

\subsection{Cell and equaliser models}
Each cell is a table-based equivalent circuit~\cite{plett2015},
\begin{equation}
v_i = V_\text{oc}(s_i) - i_i R_0(s_i), \qquad \dot s_i = -\frac{i_i}{3600\,C_n},
\label{eq:cell}
\end{equation}
with SOC $s_i$, discharge current $i_i$, capacity $C_n$ in Ah, $R_0=33$--47\,m$\Omega$ and $V_\text{oc}$
interpolated from 3.0/3.4/3.6/3.7/3.8/4.0/4.2\,V at 0/10/25/50/75/90/100\,\% SOC.

A flying capacitor $C$ alternately connected across two cells that differ by $\Delta V_k$ carries $q=C\Delta V_k$
per cycle. With switching frequency $f$, conduction fraction $D$ per position and loop resistance $R$, the pair
behaves as~\cite{seeman2008,benyaakov2012}
\begin{equation}
R_\text{eq}=\frac{1}{fC}\,\frac{1+e^{-x}}{1-e^{-x}},\qquad x=\frac{D}{fRC},
\label{eq:req}
\end{equation}
so that the average pair current is
\begin{equation}
I_k=\frac{\Delta V_k}{R_\text{eq}}, \qquad \Delta V_k=v_k-v_{k+1}.
\label{eq:ipair}
\end{equation}
Equation~\eqref{eq:req} tends to $1/(fC)$ for slow switching and to $2R/D$ for fast switching; it agreed with a
switch-level Simscape model to within 0.3\,\% from 10\,Hz to 100\,kHz (Fig.~\ref{fig:simsc}).

\subsection{Simulink implementation and results}
The model (Fig.~\ref{fig:simulink}) is generated by a MATLAB script and contains the cells with thermal ports,
bleed branches with resistor temperature, the averaged equaliser, a measurement model (quantisation, noise, gain
and offset error) and the controller as a MATLAB Function block written to translate line by line into the Arduino
loop. From 80/65/75/70\,\% SOC (107\,mV) with 2.5\,Ah cells, 22\,$\Omega$ bleed resistors and a 2.2\,mF capacitor
switched at 2\,kHz, both methods reached about 11\,mV (Table~\ref{tab:sim}, Fig.~\ref{fig:simcmp}). Passive
balancing brought every cell down to the lowest; active balancing moved the cells towards the mean and left 10\,\%
more deliverable charge. A resistor sweep gave 3.21\,h and 46\,$^\circ$C for 47\,$\Omega$, the value fitted on the
hardware.

\begin{figure*}[t]
\centering
\figslot{simulink_model.png}{0.9\textwidth}{6cm}{Simulink/Simscape model screenshot}
\caption{Simulink model of the four-cell pack and BMS (logging blocks omitted).}
\label{fig:simulink}
\end{figure*}

\begin{table}[t]
\caption{Simulated comparison (four runs per method)}
\label{tab:sim}
\centering
\small
\begin{tabular}{@{}lccc@{}}
\toprule
Metric & None & Passive & Active \\
\midrule
Spread [mV] & 107$\to$107 & 107$\to$11 & 107$\to$11 \\
Time to $\le$20\,mV [h] & -- & $1.55\pm0.02$ & $1.05\pm0.02$ \\
Energy lost [Wh] & 0 & $2.05\pm0.04$ & 0.007 \\
Peak resistor $T$ [$^\circ$C] & 25 & 69 & 25 \\
Deliverable charge [Ah] & 1.62 & 1.62 & 1.78 \\
\bottomrule
\end{tabular}
\end{table}

\begin{figure}[t]
\centering
\figslot{sim_method_comparison.png}{\columnwidth}{5cm}{Simulated comparison of the three cases}
\caption{Simulated comparison from identical initial conditions.}
\label{fig:simcmp}
\end{figure}

\begin{figure}[t]
\centering
\figslot{sim_sc_characterisation.png}{\columnwidth}{5cm}{Equivalent resistance against switching frequency}
\caption{Equivalent resistance of one capacitor stage: switch-level model and \eqref{eq:req}.}
\label{fig:simsc}
\end{figure}

\section{Controller Design}
The firmware starts in OFF (monitoring only); \texttt{p}, \texttt{a} and \texttt{o} select passive, active and off.
Every mode change switches all outputs off and begins with a rest (Fig.~\ref{fig:flow}).

\textbf{Measure-at-rest cycle.} While a cell carries current, its reading is shifted by the current through the
cell and wiring resistance. Each mode therefore alternates RUN and REST, 10\,s/5\,s for passive and 30\,s/5\,s for
active balancing, and only readings taken with every output off for at least 3\,s are used for decisions.

\textbf{Decision rules.} With $v_\text{min}=\min_i v_i$, passive balancing selects cell $i$ with hysteresis,
\begin{equation}
b_i \leftarrow
\begin{cases}
1, & v_i - v_\text{min} > \Delta_\text{on},\\
0, & v_i - v_\text{min} \le \Delta_\text{off},\\
b_i, & \text{otherwise},
\end{cases}
\label{eq:passive_rule}
\end{equation}
and active balancing runs while the spread $\sigma=\max_i v_i-\min_i v_i$ exceeds $\Delta_\text{on}$ and stops at
$\sigma\le\Delta_\text{off}$, with $\Delta_\text{on}=15$\,mV and $\Delta_\text{off}=10$\,mV (60 and 40 ADC steps).

\begin{figure}[t]
\centering
\figslot{flowchart.png}{0.92\columnwidth}{8cm}{Controller flowchart}
\caption{Controller flowchart, executed once per second.}
\label{fig:flow}
\end{figure}

\section{Calibration}
\label{sec:cal}
Each cell is measured with a multimeter at its terminals while the BMS is in OFF mode, and compared with the mean
of ten logged readings. The error $e_i=v_i^\text{BMS}-v_i^\text{DMM}$ quantifies the measurement error; a gain
$g_i=v_i^\text{DMM}/v_i^\text{BMS}$ (one point) or a gain and offset (two or more points) is entered in the
firmware and the error is measured again. Cells 3 and 4 are differences of node voltages near 12--16\,V, so 1\,\%
divider tolerance alone can cause tens of millivolts of error.

Table~\ref{tab:cal} lists the result. Before calibration the errors were $-36$, $+33$, $-48$ and $+45$\,mV
(0.9--1.2\,\%, 41\,mV RMS), consistent with 1\,\% resistors; after the gains were applied every channel read
within $\pm1$\,mV of the multimeter, which is at the level of the meter's own resolution. The runs in
Section~\ref{sec:hw} were logged with unity gains. Applying Table~\ref{tab:cal} to them changes each change of a
single cell by only about 1\,\%, but shifts the absolute spreads by $-9$ to $+83$\,mV; spread values quoted
below are therefore uncalibrated, while per-cell changes are unaffected.

\begin{table}[t]
\caption{Calibration against a multimeter (offsets $\le2$\,mV)}
\label{tab:cal}
\centering
\small
\begin{tabular}{@{}lccccc@{}}
\toprule
Cell & DMM [V] & BMS [V] & Error [mV] & Gain & After [mV] \\
\midrule
1 & 3.851 & 3.815 & $-36$ & 1.0094 & $-1$ \\
2 & 3.852 & 3.885 & $+33$ & 0.9915 & $+1$ \\
3 & 3.849 & 3.801 & $-48$ & 1.0126 & $-1$ \\
4 & 3.850 & 3.895 & $+45$ & 0.9884 & $+1$ \\
\midrule
\multicolumn{3}{@{}l}{RMS error} & 41.0 & & 1.0 \\
\bottomrule
\end{tabular}
\end{table}

\section{Estimated Balancing Time}
\label{sec:est}
Converting rested voltages to SOC with the OCV table, $s_i=V_\text{oc}^{-1}(v_i)$, passive balancing must remove
\begin{equation}
\Delta Q_i = (s_i - s_\text{min})\,C_n
\label{eq:dq}
\end{equation}
from cell $i$; with bleeding active a fraction $d_p=10/15$ of the time, the highest cell sets the time,
\begin{equation}
t_p \approx \frac{\max_i \Delta Q_i}{I_b\,d_p}.
\label{eq:tp}
\end{equation}
Charge is conserved in active balancing, so with equal capacities the cells converge to the mean SOC:
\begin{equation}
V_\infty^{\text{act}} = V_\text{oc}\!\left(\tfrac14\textstyle\sum_i s_i\right), \qquad
V_\infty^{\text{pas}} = V_\text{oc}(s_\text{min}).
\label{eq:vfinal}
\end{equation}
The active time was obtained by integrating \eqref{eq:cell} and \eqref{eq:ipair} for the three pairs with the
30\,s/5\,s cycle until $\sigma\le10$\,mV. For the pack at 4.063/4.009/3.840/3.960\,V (about 93/90/78/87\,\% SOC,
$C_n=2.6$\,Ah), \eqref{eq:dq} gives 0.395/0.324/0/0.234\,Ah, \eqref{eq:tp} gives 7.0\,h (6.7\,h by integration,
3.5\,Wh dissipated), and \eqref{eq:vfinal} gives 3.840\,V for passive and 3.962\,V for active balancing.
Table~\ref{tab:est} lists the options.

\begin{table}[t]
\caption{Estimated time to balance the pack ($\sigma\le10$\,mV)}
\label{tab:est}
\centering
\small
\begin{tabular}{@{}lrr@{}}
\toprule
Configuration & $R_\text{eq}$ [$\Omega$] & Time \\
\midrule
Passive, 47\,$\Omega$, 10\,s/5\,s & -- & 6.7\,h \\
\midrule
Relays, 470\,$\mu$F, 4.7\,$\Omega$ (built) & 2728 & 26\,700\,h \\
Relays, 4$\times$470\,$\mu$F & 682 & 6\,670\,h \\
MOSFET 100\,Hz, 470\,$\mu$F, 4.7\,$\Omega$ & 27.0 & 264\,h \\
MOSFET 1\,kHz, 470\,$\mu$F, 4.7\,$\Omega$ & 20.1 & 196\,h \\
MOSFET 1\,kHz, 4$\times$470\,$\mu$F & 5.02 & 49\,h \\
MOSFET 1\,kHz, 470\,$\mu$F, 0.4\,$\Omega$ & 2.49 & 24\,h \\
\bottomrule
\end{tabular}
\end{table}

\section{Hardware Results}
\label{sec:hw}
Table~\ref{tab:runs} lists the recorded runs. Readings were selected by three fixed rules, applied identically to
every cell and run: (i) only rested rows (all outputs off for at least 3\,s, status OK); (ii) a rested reading more
than 3\,mV from the median of its seven neighbours is discarded, since bleeding moves a cell by only about
0.2\,mV per 15\,s cycle; (iii) a cell whose noise exceeds 1\,mV or whose 1\,min rate of change exceeds
10\,mV/min is marked contact-affected and not used for conclusions. Real cells changed by at most 3.3\,mV/min
(polarisation when bleeding starts), whereas the intermittent contact on cells 3 and 4 produced 19--44\,mV/min.
Readings taken while outputs are on are shifted and are used only to characterise that shift.

\begin{table}[t]
\caption{Hardware runs}
\label{tab:runs}
\centering
\small
\begin{tabular}{@{}clcl@{}}
\toprule
Run & Mode, capacitor & Duration & Events \\
\midrule
1 & active, 470\,$\mu$F & 93.6\,min & open sense wire (FAULT) \\
2 & active, 1\,$\mu$F & 2.7\,min & none \\
3 & active, 470\,$\mu$F & 53.5\,min & cell-4 contact (WARN) \\
4 & passive, 47\,$\Omega$ & 35.6\,min & cell-4 contact (WARN) \\
\bottomrule
\end{tabular}
\end{table}

\subsection{Fault detection: open sense wire}
In run~1, 93 rows between 6 and 28\,min were flagged FAULT (Fig.~\ref{fig:hwv}). Cell~4 read $-8.192$\,V,
the negative full scale ($-2.048\,\text{V}\times4$) of $A2-A3$, which identifies an open lead on the pack-top
node, and one temperature sensor failed. After the lead was re-seated no fault occurred. Later intermittent dips
of cell~4 (to 3.06\,V in run~3 and 3.71\,V in run~4) were caught by the jump warning, and in run~3 the
low-voltage check ended the RUN phase within one second.

\begin{figure*}[t]
\centering
\figslot{hw_run1_voltages.png}{0.85\textwidth}{6cm}{Hardware run 1 cell voltages}
\caption{Run 1: cell voltages; the shaded interval is the open sense wire.}
\label{fig:hwv}
\end{figure*}

\subsection{Passive balancing}
Run~4 lasted 35.6\,min. Following \eqref{eq:passive_rule}, the controller bled cells 1 and 2 almost
throughout and cells 3 and 4 only while they were more than 15\,mV above the lowest
(Fig.~\ref{fig:hwpasout}). Cells 1 and 2 passed all selection rules (noise 0.30 and 0.25\,mV); cells 3 and 4 were
contact-affected. With $I_b=v_i/47\,\Omega$, the log gives 34.0 and 29.8\,mAh bled from cells 1 and 2
(Fig.~\ref{fig:hwcharge}, Table~\ref{tab:pas}) and 0.311\,Wh dissipated in total. Their rested voltages fell
steadily by 14.1 and 23.1\,mV (Fig.~\ref{fig:hwpas}). Through the OCV table the removed charge corresponds to 26
and 15\,mV, 41\,mV together, against 37\,mV measured; the split between the two cells differs because the OCV
table is generic and cell~1 was still recovering from an earlier drain (it rose by 17.6\,mV in run~3). Cell
temperatures stayed between 16.6 and 17.3\,$^\circ$C; the resistor temperature was not measured. Balancing was not
completed: 36\,min is short compared with the 6.7\,h estimate.

\begin{table}[t]
\caption{Passive run (run 4, 35.6 min)}
\label{tab:pas}
\centering
\small
\begin{tabular}{@{}lccccc@{}}
\toprule
Cell & On [min] & $\Delta Q$ [mAh] & $E$ [Wh] & $\Delta V$ [mV] & Exp. [mV] \\
\midrule
1 & 23.7 & 34.0 & 0.137 & $-14.1$ & $-26$ \\
2 & 21.5 & 29.8 & 0.116 & $-23.1$ & $-15$ \\
3 & 2.0 & 2.8 & 0.011 & \multicolumn{2}{c}{contact-affected} \\
4 & 8.5 & 11.8 & 0.046 & \multicolumn{2}{c}{contact-affected} \\
\midrule
Total & & 78.3 & 0.311 & & \\
\bottomrule
\end{tabular}
\end{table}

\begin{figure}[t]
\centering
\figslot{hw_run4_passive_outputs.png}{\columnwidth}{4.5cm}{Passive run: mode and bleed outputs}
\caption{Run 4: bleed relays selected by the controller.}
\label{fig:hwpasout}
\end{figure}

\begin{figure}[t]
\centering
\figslot{hw_run4_clean_change.png}{\columnwidth}{4.5cm}{Passive run: rested change of each cell, selected readings}
\caption{Run 4: change of the rested cell voltages after the selection rules. Cells 1 and 2 fall steadily;
cells 3 and 4 (faded) are contact-affected.}
\label{fig:hwpas}
\end{figure}

\begin{figure}[t]
\centering
\figslot{hw_run4_clean_charge.png}{\columnwidth}{4.2cm}{Passive run: cumulative charge bled}
\caption{Run 4: cumulative charge bled from each cell.}
\label{fig:hwcharge}
\end{figure}

\subsection{Active balancing}
Over the stable hour of run~1 the spread went from 223.3 to 225.5\,mV, and over run~3 (93 cycles, 2652 shuttles)
from 185.3 to 187.6\,mV; the reliable cells 1 and 2 rose by 17.6 and 4.2\,mV instead of moving towards
the lower cells (Fig.~\ref{fig:hwr3}). Applying \eqref{eq:req} with $C=470\,\mu$F and $R=4.8\,\Omega$
gives $R_\text{eq}=2.7$\,k$\Omega$ for run~1 ($f=0.78$\,Hz) and 2.1\,k$\Omega$ for run~3 ($f=1$\,Hz), and
\eqref{eq:ipair} gives 20--60\,$\mu$A per pair, about 0.002\,\% SOC per hour, below one ADC step per hour. The
capacitors did switch: a multimeter showed each between its two cell voltages, and in run~1 cells 1 and 2 showed a
$-2.3/+1.1$\,mV step at every transition from RUN to REST. Replacing the capacitors with 1\,$\mu$F (run~2) removed
this step ($<0.05$\,mV), confirming its origin. Run~2 also serves as the monitoring-only reference, since its
transfer current is about 0.1\,$\mu$A.

\begin{figure}[t]
\centering
\figslot{hw_run3_clean_change.png}{\columnwidth}{4.5cm}{Run 3: rested change of each cell, selected readings}
\caption{Run 3: change of the rested cell voltages during active balancing after the selection rules.}
\label{fig:hwr3}
\end{figure}

\subsection{Measurement shifts during RUN}
While outputs are on, the cell-1 reading is shifted by 226\,mV with the capacitor relays (run~3) and by 351\,mV
with the bleed relays (run~4); cells 2 and 4 shift by 132 and 73\,mV during bleeding (Fig.~\ref{fig:hwraw}).
These shifts are far larger than the cell drop $I_bR_0\approx3$\,mV and appear as soon as the relays switch, so
they are attributed to relay-coil and bleed currents flowing in wiring shared with the measurement, in particular
the ground of the single-ended cell-1 channel. Because decisions use only rested readings, the controller is
unaffected.

\begin{figure}[t]
\centering
\figslot{hw_run3_raw_voltages.png}{\columnwidth}{4.5cm}{Run 3 every row, showing the RUN shift}
\caption{Run 3, every row: cell~1 reads 226\,mV low while the relays shuttle.}
\label{fig:hwraw}
\end{figure}

\subsection{Comparison}
Table~\ref{tab:cmp} compares the cases on hardware.

\begin{table}[t]
\caption{Comparison on hardware}
\label{tab:cmp}
\centering
\small
\begin{tabular}{@{}lccc@{}}
\toprule
Metric & Reference & Passive & Active \\
 & (run 2) & (run 4) & (run 3) \\
\midrule
Duration [min] & 2.7 & 35.6 & 53.5 \\
$\Delta V_1$ / $\Delta V_2$ [mV] & 0 / 0 & $-14.1$ / $-23.1$ & $+17.6^{\ast}$ / $+4.2$ \\
Charge moved [mAh] & -- & 78.3 & $\approx0.05$ \\
Energy [Wh] & 0 & 0.311 & $<10^{-4}$ \\
Peak cell $T$ [$^\circ$C] & -- & 17.3 & 17.9 \\
Est. time to balance & -- & 6.7\,h & 26\,700\,h \\
\bottomrule
\multicolumn{4}{@{}l}{\footnotesize $^{\ast}$likely recovery of cell 1, which had been drained before the run}
\end{tabular}
\end{table}

\section{Discussion}
\label{sec:disc}
\textbf{Passive balancing works as designed.} The controller chose the right cells, the bleed current matched
\eqref{eq:bleed}, and the combined rested fall of cells 1 and 2 (37\,mV) matched the 41\,mV expected from the
removed charge. Completing the balance needs a
run of about 7\,h; a 22\,$\Omega$ resistor would halve this at about 0.75\,W per resistor.

\textbf{Switching rate sets active performance.} By \eqref{eq:req}, a capacitor stage moves charge in proportion
to $fC$ until the loop resistance dominates. The simulated stage used $fC=4.4$\,F/s; the relay stage has
$fC\approx4.7\times10^{-4}$\,F/s, about 9\,000 times less. Replacing each relay by a logic-level MOSFET half-bridge
(an N/P pair across each cell, gates driven through digital isolators powered by that cell, Timer1 generating both
phases with dead time) at 1\,kHz brings the time to 24--196\,h with the existing capacitors (Table~\ref{tab:est}).
A fuse in every cell tap is then required, because simultaneous conduction shorts a cell. The 4.7\,$\Omega$
resistors cap such a stage at about 8\,mA on a 168\,mV pair, so a lower loop resistance with a precharge step is
needed for faster balancing.

\textbf{Divider loading.} Each 44\,k$\Omega$ divider runs from its node to ground, so cell~$i$ supplies the
dividers $j\ge i$,
\begin{equation}
I_{\text{drain},i}=\sum_{j\ge i}\frac{V_{\text{node},j}}{R_\text{top}+R_\text{bot}},
\label{eq:drain}
\end{equation}
which is 0.90/0.81/0.63/0.36\,mA for cells 1--4: the difference exceeds the transfer of the relay stage by an
order of magnitude.

\textbf{Wiring.} The intermittent cell-4 contact and the RUN shifts both come from wiring. Separate sense and
current-carrying wires (Kelvin connections), a separate supply and ground return for the relay coils, and
soldered connections would remove both.

\section{Limitations and Uncertainty}
\begin{itemize}
\item The runs were logged before the calibration gains were applied (Section~\ref{sec:cal}); absolute spreads
are uncertain by up to about 80\,mV, while changes of one cell over a run are resolved to about 0.5\,mV and are
affected by only 1\,\%. Calibration used a single point near 3.85\,V per cell.
\item No capacitor test emulator was used before the cells were connected.
\item The passive run lasted 36\,min, far shorter than the 6.7\,h needed to balance; no repeat runs and no
dedicated no-balancing run of equal length were recorded, so repeatability is not quantified.
\item The intermittent cell-4 / node-A3 contact affects cells 3 and 4 in runs 1, 3 and 4; their readings were
excluded by the selection rules, so conclusions rest on cells 1 and 2.
\item The OCV table is generic, so SOC values and time estimates are approximate; bleed currents are inferred from
voltage and the nominal 47\,$\Omega$, not measured directly; resistor temperature was not measured.
\item The active stage was built only with relays; the MOSFET stage is designed but not built.
\end{itemize}

\section{Conclusion}
A single Arduino controller provides monitoring, passive and active balancing, alerts, latched faults and a manual
stop; its fault and warning logic caught a real open sense wire and later intermittent contacts, and calibration
against a multimeter brought the measurement error from up to 48\,mV to $\pm1$\,mV. Simulation
predicted that both methods balance the pack, active balancing faster and nearly losslessly. On hardware, passive
balancing behaved as designed: the controller selected the high cells, and the measured voltage fall of the two highest cells matched the 64\,mAh removed at about 85\,mA, with
0.31\,Wh dissipated in 36\,min and no measurable heating. The
relay-switched 470\,$\mu$F ladder moved only microamperes because it switches about once per second, in agreement
with the averaged model. The same model and the measured voltages give 6.7\,h for full passive balancing and
24--196\,h for a 1\,kHz MOSFET version of the active stage, which is the main improvement needed for active
balancing to be competitive.

\section*{Code and Data Availability}
The Arduino firmware, MATLAB/Simulink models, Python analysis scripts, raw serial logs, calibration data,
schematics and all graphs are available at \url{https://github.com/highdreameater/EE-302-Control-sys-proj}.

\begin{thebibliography}{10}

\bibitem{daowd2011}
M.~Daowd, N.~Omar, P.~Van Den Bossche, and J.~Van Mierlo, ``Passive and active battery balancing comparison based
on MATLAB simulation,'' in \emph{Proc. IEEE Vehicle Power Propulsion Conf. (VPPC)}, Chicago, IL, 2011, pp.~1--7.

\bibitem{seeman2008}
M.~D. Seeman and S.~R. Sanders, ``Analysis and optimization of switched-capacitor DC--DC converters,''
\emph{IEEE Trans. Power Electron.}, vol.~23, no.~2, pp.~841--851, Mar. 2008.

\bibitem{pascual1997}
C.~Pascual and P.~T. Krein, ``Switched capacitor system for automatic series battery equalization,'' in
\emph{Proc. IEEE Appl. Power Electron. Conf. (APEC)}, Atlanta, GA, 1997, vol.~2, pp.~848--854.

\bibitem{gallardo2014}
J.~Gallardo-Lozano, E.~Romero-Cadaval, M.~I. Milan\'es-Montero, and M.~A. Guerrero-Mart\'inez, ``Battery
equalization active methods,'' \emph{J. Power Sources}, vol.~246, pp.~934--949, 2014.

\bibitem{benyaakov2012}
S.~Ben-Yaakov, ``Behavioral average modeling and equivalent circuit simulation of switched capacitors
converters,'' \emph{IEEE Trans. Power Electron.}, vol.~27, no.~2, pp.~632--636, Feb. 2012.

\bibitem{plett2015}
G.~L. Plett, \emph{Battery Management Systems}, vols. I and II. Norwood, MA: Artech House, 2015.

\bibitem{ads1115}
Texas Instruments, \emph{ADS111x Ultra-Small, Low-Power, I$^2$C-Compatible, 16-Bit ADCs}, datasheet.

\bibitem{lim2025}
W.~C. Lim, L.~Siek, and E.~L. Tan, ``Adaptive hybrid switched-capacitor cell balancing for 4-cell Li-ion battery
pack with a study of pulse-frequency modulation control,'' \emph{J. Low Power Electron. Appl.}, vol.~15, no.~4,
art.~61, 2025.

\end{thebibliography}

\end{document}
